\chapter{Production Strategy and Organization Plan}
\label{ch:organization}

\section{The process of obtaining the service}

The service is delivered entirely online, and its production is therefore not
manufacturing but operation: keeping a platform available, keeping the matching
engine accurate, and arbitrating the human steps that software cannot decide.
Figure~\ref{fig:service-process} shows how the three actors converge.

\begin{figure}[htbp]
\centering
\begin{tikzpicture}[x=1cm,y=1cm]

  \node[boxnode, text width=3.4cm] (loser)  at (-4.6,6.0) {\textbf{Person who lost}\\an object};
  \node[boxnode, text width=3.4cm] (finder) at ( 0.0,6.0) {\textbf{Person who found}\\an object};
  \node[boxnode, text width=3.4cm] (inst)   at ( 4.6,6.0) {\textbf{Partner institution}\\lost-property desk};

  \node[storenode, text width=3.4cm] (rlost)  at (-4.6,4.3) {Files a \emph{lost} report\\{\scriptsize description, photos, wilaya, date}};
  \node[storenode, text width=3.4cm] (rfound) at ( 0.0,4.3) {Files a \emph{found} report\\{\scriptsize + verification questions}};
  \node[storenode, text width=3.4cm] (rinst)  at ( 4.6,4.3) {Publishes its register\\{\scriptsize as found reports}};

  \node[boxnode, text width=12.6cm] (engine) at (0,2.4)
    {\textbf{Platform} --- embed $\rightarrow$ retrieve $\rightarrow$ score $\rightarrow$ rank\\
     {\scriptsize automatic, bidirectional, continuous}};

  \node[storenode, text width=12.6cm] (sugg) at (0,0.9)
    {Ranked suggestions with confidence and explanations; both owners notified above threshold;\\
     {\scriptsize first unlock free, then one credit per suggestion}};

  \node[storenode, text width=12.6cm] (claim) at (0,-0.6)
    {Claim submitted --- the claimant answers the verification questions set by the reporter};

  \node[storenode, text width=12.6cm] (approve) at (0,-2.1)
    {\textbf{The reporter} reviews the answers and approves or rejects};

  \node[boxnode, text width=12.6cm] (done) at (0,-3.6)
    {Contact details exchanged $\rightarrow$ handover $\rightarrow$ both reports closed as \emph{recovered}\\
     {\scriptsize the outcome is stored as labelled feedback for the matching engine}};

  \draw[flowarrow] (loser)  -- (rlost);
  \draw[flowarrow] (finder) -- (rfound);
  \draw[flowarrow] (inst)   -- (rinst);
  \draw[flowarrow] (rlost.south)  |- ($(engine.north)+(0,0.35)$) -- (engine.north);
  \draw[flowarrow] (rfound.south) -- ($(engine.north)+(0,0.35)$) -- (engine.north);
  \draw[flowarrow] (rinst.south)  |- ($(engine.north)+(0,0.35)$) -- (engine.north);
  \draw[flowarrow] (engine)  -- (sugg);
  \draw[flowarrow] (sugg)    -- (claim);
  \draw[flowarrow] (claim)   -- (approve);
  \draw[flowarrow] (approve) -- (done);

\end{tikzpicture}
\caption{Delivery of the service: three entry points, one automatic comparison,
and a human decision before any identity is disclosed.}
\label{fig:service-process}
\end{figure}

\subsection{For the person who lost an object}

\begin{enumerate}[leftmargin=1.6em,itemsep=3pt]
  \item Creates an account and files a report describing the object, in the
        language of their choice, with photographs if any exist.
  \item Receives ranked suggestions with confidence and explanations as soon as
        matching runs, and a notification whenever a strong new match appears.
        The first suggestion they open is free; each further one costs a
        credit.
  \item Submits a claim on a promising found report, answering the verification
        questions set by the finder.
  \item On approval, receives the finder's contact details and arranges the
        handover.
  \item Closes the report as recovered, which supplies the labelled feedback the
        matching engine learns from.
\end{enumerate}

Nothing in this sequence requires the user to search. That is the point: the
comparison work that social-media groups leave to human attention is performed
by the platform, continuously, for as long as the report stays open.

\subsection{For the person who found an object}

The path mirrors the one above, with one difference that determines whether the
platform works at all. Found reports are the scarce side of this market: people
hesitate to publish that they are holding something valuable, because doing so
in a public group means exposing themselves to every opportunist who reads it.

The platform answers that objection structurally rather than with reassurance:

\begin{itemize}[leftmargin=1.4em,itemsep=3pt]
  \item \textbf{No contact details are ever published.} A found report shows the
        object, not the finder.
  \item \textbf{The finder sets the verification questions} when publishing
        --- \dq{which side is it, left or right?}, \dq{what is on the
        lock screen?} --- and chooses questions only the true owner can answer.
  \item \textbf{The finder alone decides} who is genuine. A high confidence
        score is a suggestion to the finder, never an instruction, and never an
        authorisation to disclose anything.
  \item \textbf{Claims cannot be used to harass.} Only one pending claim per
        person per item is permitted, so a rejected claimant cannot flood the
        finder with retries.
\end{itemize}

A finder is therefore never more exposed than they choose to be, which is the
minimum condition for the found pool to grow at all.

\subsection{For partner institutions}

Universities, transport operators, municipalities and commercial centres already
accumulate lost property, and already keep registers of it. What they lack is
any way to make those registers findable. A citizen who lost something must
guess which desk to visit, travel there, and ask during opening hours ---
repeated once per institution.

The proposed arrangement is simple: the institution receives an account, its
staff publish found items into the shared pool as they are handed in, and claims
arrive through the same verification workflow that protects individual finders.
What the institution gains:

\begin{itemize}[leftmargin=1.4em,itemsep=3pt]
  \item A searchable register replacing a paper one, at no licence cost.
  \item Fewer speculative counter visits, because citizens can check before
        travelling.
  \item Automatic matching against loss reports, instead of waiting for the
        owner to think of asking.
  \item A record of restitutions, which is evidence of service quality.
\end{itemize}

For the platform, these partnerships are not merely a channel --- they are the
answer to the cold-start problem identified in Chapter~\ref{ch:market}. An
institutional register seeds the found pool on day one, so that the first
individual user to arrive after a loss finds something to match against.

\section{Procurement}

The stack is built entirely on open-source components --- PostgreSQL, Redis,
FastAPI, Next.js --- and on pre-trained models used under permissive licences
and run locally. Procurement is therefore a handful of monthly subscriptions;
there are no software licences to buy, no equipment to purchase, and,
critically, \textbf{no per-request charge to any external AI provider}. Every
embedding is computed on the project's own CPU, which is what keeps reporting
free and the price of an unlock low. Table~\ref{tab:procurement} lists what is
procured; the prices are justified in Chapter~\ref{ch:finance}.

\begin{table}[htbp]
\centering
\caption{Procurement requirements}
\label{tab:procurement}
\small
\begin{tabularx}{\textwidth}{@{}L{3.2cm} X L{3.9cm}@{}}
\toprule
\textbf{Item} & \textbf{Why it is needed} & \textbf{Offer chosen} \\
\midrule
One virtual server
  & Runs all five services: database, queue, API, worker and frontend. Sized
    on the measured 1.2\,GB of memory; no GPU needed
  & Oracle Always Free for the demonstration; Hetzner CX23, 5.49\,EUR a
    month, for the pilot~\cite{hetzner2026} \\
\addlinespace
Public IPv4 and backups
  & A reachable address, and daily copies of the database and photographs
  & Hetzner, 0.50\,EUR a month plus 20\,\% of the server
    price~\cite{cloudtally2026} \\
\addlinespace
Domain name
  & The platform's public identity
  & \code{.dz} domain, 1\,500\,DA a year~\cite{hostarts2026} \\
\addlinespace
HTTPS certificates
  & Encrypted connections
  & Let's Encrypt through Caddy, free~\cite{caddy2026} \\
\addlinespace
Transactional e-mail
  & Account and password e-mails as the user base grows
  & Brevo free plan, 300 e-mails a day~\cite{brevo2026} \\
\addlinespace
Payment gateway
  & Selling credit packs in dinars, by EDAHABIA or CIB card
  & Chargily Pay: free for 300 transactions, then 1.5\,\% with a 15\,DZD
    minimum~\cite{chargily2026} \\
\bottomrule
\end{tabularx}
\end{table}

Two architectural decisions keep this list short and its total low.
Photographs are re-encoded to WebP at a bounded resolution before storage, and
those of the evaluation corpus average 63\,KB, so storage and bandwidth are not
procured separately. And both models run on the CPU within 872\,MB of memory,
so a single 4\,GB server hosts the whole platform, database included, instead
of a managed database, an object store and a GPU instance billed separately.

\section{Workforce}

The team required is small, because the product's marginal cost is
computational rather than human. What does scale with volume is moderation:
every disputed claim is arbitrated by a person, and claim volume grows with
reports. The administration console of Chapter~\ref{ch:prototype} lets the
founders do that work themselves during the pilot.

\begin{table}[htbp]
\centering
\caption{Who does what, by phase (the phases of Chapter~\ref{ch:finance})}
\label{tab:headcount}
\small
\begin{tabularx}{\textwidth}{@{}X X X X@{}}
\toprule
\textbf{Role} & \textbf{Phase 0 --- demonstration} & \textbf{Phase 1 --- pilot} & \textbf{Phase 2 --- company} \\
\midrule
Development and operation & Founding team, unpaid & Founding team, unpaid & Two salaried developers \\
\addlinespace
Moderation and support    & Founding team & Founding team, through the console & The developers, then a moderator \\
\addlinespace
Institutional partnerships & Founding team & Founding team & Founding team \\
\addlinespace
Accounting and declarations & --- & Founders (simplified regime) & External accountant \\
\midrule
\textbf{Salaried staff}   & \textbf{0} & \textbf{0} & \textbf{2} \\
\bottomrule
\end{tabularx}
\end{table}

No salary is paid before the pilot has shown recovery volume and a measured
conversion rate; the first two salaries are the ones costed in
Chapter~\ref{ch:finance}. The next hire is a moderator, once claim volume
exceeds what the developers can arbitrate while still building.

\section{Key partners}

\begin{itemize}[leftmargin=1.4em,itemsep=5pt]
  \item \textbf{Universities.} The pilot ground. Akli Mohand Oulhadj University
        of Bouira first, then the wider network of more than 130
        institutions~\cite{mesrs2025}. They contribute a bounded population with
        high object turnover and an existing lost-property desk; they receive a
        searchable register and fewer counter visits.

  \item \textbf{Public transport operators} --- metro, tramway, rail and bus
        networks --- where loss events concentrate most heavily. They contribute
        volume on the scarce found side; they receive automated matching against
        national loss reports.

  \item \textbf{Airports and stations}, which already operate formal
        lost-property procedures and legal retention periods, and for whom an
        online register reduces counter load.

  \item \textbf{Municipalities and communal services}, extending coverage beyond
        institutional spaces into the street.

  \item \textbf{Commercial centres and large retailers}, which hold objects
        found on their premises and gain a visible customer service.

  \item \textbf{The university incubator and Algeria Venture}, the national
        public accelerator supporting labelled startups~\cite{startupdz}, for
        mentoring, credibility and access to the financing mechanisms of
        Chapter~\ref{ch:finance}.

  \item \textbf{ANADE}, whose interest-free financing schemes for young ventures
        reach up to 10 million dinars~\cite{anade2025}.

  \item \textbf{Hosting providers} --- Oracle's free tier for the
        demonstration, Hetzner for the pilot --- whose offers let the platform
        run for under 2\,000~DZD a month.

  \item \textbf{Chargily}, the payment gateway through which credit packs are
        sold in dinars by national bank card.
\end{itemize}
